% --- （ここにあなたが用意されたパッケージ群をそのまま貼りました） ---
\usepackage{comment}    % 複数行コメントを可能にする
\usepackage{float}      % 図表の配置を制御
\usepackage{color}      % 色を扱う
\usepackage{multicol}   % 複数列レイアウトをサポート
\usepackage[dvipdfmx]{pict2e} % 高品質な描画命令（日本語PDF対応）
\usepackage{wrapfig}    % 本文中に図をラップさせる
\usepackage{graphicx}   % 画像挿入用
\usepackage{url}        % URLの適切な表示
\usepackage{underscore} % アンダースコアを文字列として扱う
\usepackage{colortbl}   % 表のセルに色をつける
\usepackage{tabularx}   % 自動列幅調整可能な表の作成
\usepackage{fancyhdr}   % ヘッダー・フッターのカスタマイズ
\usepackage{ulem}       % 下線や取り消し線の追加
\usepackage{cite}       % 文献の引用を整理
\usepackage{amsmath,amssymb,amsfonts} % 数式関連の基本パッケージ群
\usepackage{algorithmic} % アルゴリズムの記述
\usepackage{textcomp}   % 特殊記号をサポート
\usepackage{xcolor}     % 拡張された色設定を使用可能にする
\usepackage[ipaex]{pxchfon} % 日本語フォント（IPAフォント）を指定
\usepackage{amsthm}     % 定理環境をサポート
\usepackage{dsfont}     % ブラックボードボールド体などをサポート
\usepackage{mathtools}  % 数式関連の拡張
\usepackage{siunitx}    % 科学技術単位を統一した形式で表記
\usepackage[thicklines]{cancel} % 数式内での取り消し線
\renewcommand{\CancelColor}{\color[rgb]{0.7,0,0}}
\usepackage{nccmath}    % 数式のスペースを調整
\usepackage{enumitem}  % 箇条書きのカスタマイズ

% hyperrefと色の設定
\usepackage[
    colorlinks=true,    
    linkcolor=blue,     % 内部リンクは青色
    citecolor=green     % 参考文献は緑色
]{hyperref}

\renewcommand{\bibname}{参考文献}

\usepackage{tcolorbox} % 箱を作成
\tcbuselibrary{breakable, skins, theorems}
\usepackage{tikz}
\usetikzlibrary{calc,decorations.pathmorphing,arrows.meta,shadows,shapes.geometric,arrows}

\usepackage{ifthen}
\usepackage{pgfplots}
\pgfplotsset{compat=1.17}

%%%%%%mathtools%%%%%
\allowdisplaybreaks % align pagebreak
%\mathtoolsset{showonlyrefs=true} % 被参照数式のみ数式番号割り振り
\numberwithin{equation}{section} % 数式番号をセクションごとにリセット

\makeatletter
\@addtoreset{equation}{section}
\makeatother

%---------- 追加: 定理環境・スタイル ----------%

%%%%%%theorem environments%%%%%
\newtheoremstyle{mystyle}%   % スタイル名
    {}%                      % 上部スペース
    {}%                      % 下部スペース
    {\normalfont}%          % 本文フォント
    {}%                      % インデント量
    {\bf}%                  % 見出しフォント
    {.}%                      % 見出し後の句読点
    {\newline}%                     % 見出し後のスペース
    {\underline{\thmname{#1}\thmnumber{#2}\thmnote{（#3）}}}%
    % 見出しの書式 (can be left empty, meaning `normal')
\theoremstyle{mystyle} % スタイルの適用

\newtheorem{theorem}{定理}[section]
\newtheorem{definition}{定義}[section]
\newtheorem{proposition}[definition]{命題}
\newtheorem{corollary}[theorem]{系}
\newtheorem{example}{例}[section]

\theoremstyle{remark}
\newtheorem{remark}{注意}[section]

%% === 証明環境の外観 ===
\makeatletter
\renewcommand{\qedsymbol}{$\blacksquare$}
\renewenvironment{proof}[1][\proofname]{\par
    \pushQED{\qed}%
    \normalfont \topsep6\p@\@plus6\p@\relax
    \trivlist
    \item[\hskip\labelsep
        \itshape
    \textbf{#1}]\ignorespaces
}{%
    \popQED\endtrivlist\@endpefalse
}
\makeatother

% ① カウンターの定義
\newcounter{defcounter}[section] % sectionごとにリセットされるカウンターを作成
\renewcommand{\thedefcounter}{\thesection.\arabic{defcounter}} % 表示形式を「1.1」のようにする

% ② 新しい環境「mydefinition」の定義
% #1 には「無作為標本」などのタイトルが入る
\newenvironment{mydefinition}[1]{
  \refstepcounter{defcounter} % カウンターを1つ増やし、参照可能にする
  \begin{definitionbox}{\underline{定義 \thedefcounter: \quad #1} } % 自動で番号を付けてdefinitionboxを開始
}{
  \end{definitionbox} % definitionboxを終了
}
% --------------------------
% ---- Definition Palette (緑系・好みで調整) ----
\definecolor{DfGreen}{RGB}{32,138,87}    % アクセント緑（落ち着きめ）
\colorlet{DfGreenLine}{DfGreen!80!black} % 線用（やや落とした緑）
\colorlet{DfGreenLite}{DfGreen!6!white}  % ごく薄い緑の塗り
\colorlet{DfGreenHair}{DfGreen!30!black} % 補助線（必要なら）
% -----------------------------------------------

% ---- definitionbox（theorembox と同じ骨格）----
\newtcolorbox{definitionbox}[2][]{
  enhanced, breakable,
  colback=DfGreenLite, colframe=white, frame hidden, boxrule=0pt,
  borderline west={1pt}{0pt}{DfGreenLine}, % 左キーラインのみ
  top=2mm, bottom=2mm, left=1.6mm, right=1.6mm,
  before skip=6pt, after skip=10pt,
  fonttitle=\bfseries, coltitle=black, title=#2,
  #1
}

% --- 共通カウンター（セクションごとにリセット） ---
\newcounter{theocounter}[section]
\renewcommand{\thetheocounter}{\thesection.\arabic{theocounter}}

% ---- Blue Palette（既存と同じ） ----
\definecolor{ThBlue}{RGB}{36,99,190}      % アクセント青（中間）
\colorlet{ThBlueLine}{ThBlue!100!black}   % 線用
\colorlet{ThBlueLite}{ThBlue!6!white}     % ごく薄い塗り
\colorlet{ThBlueHair}{ThBlue!30!black}    % ヘアライン（未使用でもOK）
% ------------------------------------

% ---- 共通ボックス骨格（既存と同じ見た目） ----
\newtcolorbox{theorembox}[2][]{
  enhanced, breakable,
  colback=ThBlueLite, colframe=white, frame hidden, boxrule=0pt,
  borderline west={1pt}{0pt}{ThBlueLine},
  top=2mm, bottom=2mm, left=1.6mm, right=1.6mm,
  before skip=6pt, after skip=10pt,
  fonttitle=\bfseries, coltitle=black, title=#2,
  #1
}
% ------------------------------------

% ===== 定理 =====
% 使い方: \begin{mytheorem}[<label>]{<定理の名称>} ... \end{mytheorem}
\newenvironment{mytheorem}[2][]
{
  \refstepcounter{theocounter}
  \ifthenelse{\equal{#1}{}}
    {\begin{theorembox}{\underline{定理 \thetheocounter: \quad #2}}}
    {\begin{theorembox}[label={#1}]{\underline{定理 \thetheocounter: \quad #2}}}
}
{\end{theorembox}}

% ===== 命題 =====
% 使い方: \begin{myproposition}[<label>]{<命題の名称>} ... \end{myproposition}
\newenvironment{myproposition}[2][]
{
  \refstepcounter{theocounter}
  \ifthenelse{\equal{#1}{}}
    {\begin{theorembox}{\underline{命題 \thetheocounter: \quad #2}}}
    {\begin{theorembox}[label={#1}]{\underline{命題 \thetheocounter: \quad #2}}}
}
{\end{theorembox}}

% ===== 補題 =====
% 使い方: \begin{mylemma}[<label>]{<補題の名称>} ... \end{mylemma}
\newenvironment{mylemma}[2][]
{
  \refstepcounter{theocounter}
  \ifthenelse{\equal{#1}{}}
    {\begin{theorembox}{\underline{補題 \thetheocounter: \quad #2}}}
    {\begin{theorembox}[label={#1}]{\underline{補題 \thetheocounter: \quad #2}}}
}
{\end{theorembox}}



% =========================
%  例 (Example) 環境一式
% =========================

% ① 例用のカウンター（セクションごとにリセット）
\newcounter{excounter}[section]
\renewcommand{\theexcounter}{\thesection.\arabic{excounter}}

% ---- Palette (好みで調整) ----
\definecolor{ExOrange}{RGB}{209,118,30}   % アクセント橙（落ち着きめ）
\colorlet{ExOrangeLine}{ExOrange!95!black}
\colorlet{ExOrangeLite}{ExOrange!6!white}
\colorlet{ExOrangeHair}{ExOrange!30!black}
% --------------------------------

% ② tcolorbox の骨格（definition/theorem と同様）
\newtcolorbox{examplebox}[2][]{
  enhanced, breakable,
  colback=ExOrangeLite, colframe=white, frame hidden, boxrule=0pt,
  borderline west={1pt}{0pt}{ExOrangeLine},
  top=2mm, bottom=2mm, left=1.6mm, right=1.6mm,
  before skip=6pt, after skip=10pt,
  fonttitle=\bfseries, coltitle=black, title=#2,
  #1
}

% ③ 新しい環境「myexample」の定義
% 使い方: \begin{myexample}[<ラベル>]{<例の名称>}
% 例: \begin{myexample}[ex:lln]{大数の法則の直感}
\newenvironment{myexample}[2][]
{
  \refstepcounter{excounter}% カウンターを進めて参照可能に
  \ifthenelse{\equal{#1}{}}
  { % ラベルなし
    \begin{examplebox}{\underline{例 \theexcounter: \quad #2}}
  }
  { % ラベルあり
    \begin{examplebox}[label={#1}]{\underline{例 \theexcounter: \quad #2}}
  }
}
{
  \end{examplebox}
}


\newtcolorbox{bluebox}{
  breakable,
  colback=blue!5!white,
  colframe=blue!75!black,
  left=1mm, right=1mm, top=1mm, bottom=1mm,
  boxrule=0.5pt,
  arc=2pt
}

\newtcolorbox{remarkbox}{
  breakable,
  colback=yellow!5,
  colframe=yellow!60!black,
  left=1mm, right=1mm, top=1mm, bottom=1mm,
  boxrule=0.5pt,
  arc=2pt
}
\newtcolorbox{greenbox}{
  breakable,
  colback=green!5,
  colframe=green!50!black,
  left=1mm, right=1mm, top=1mm, bottom=1mm,
  boxrule=0.5pt,
  arc=2pt
}

% =========================
%  Yellow Note/Attention Boxes（カウンター無し）
% =========================

% ---- Yellow Palette ----
\definecolor{NtYellow}{RGB}{230,170,0}     % アクセント黄（アンバー寄り）
\colorlet{NtYellowLine}{NtYellow!90!black} % 左ライン用
\colorlet{NtYellowLite}{NtYellow!7!white}  % 薄い塗り
% ------------------------

% ---- 共通ボックス骨格 ----
\newtcolorbox{notebox}[2][]{
  enhanced, breakable,
  colback=NtYellowLite, colframe=white, frame hidden, boxrule=0pt,
  borderline west={1pt}{0pt}{NtYellowLine}, % 左キーライン
  top=2mm, bottom=2mm, left=1.6mm, right=1.6mm,
  before skip=6pt, after skip=10pt,
  fonttitle=\bfseries, coltitle=black, title=#2,
  #1
}

% =========================
%  環境：補足／メモ／注意
%  使い方：
%   \begin{myproofnote} ... \end{myproofnote}
%   \begin{mynote} ... \end{mynote}
%   \begin{myattention} ... \end{myattention}
%  任意タイトルにしたいとき：
%   \begin{mynote}[計算メモ] ... \end{mynote}
% =========================

% 補足（デフォルト見出し：補足）
\newenvironment{myproofnote}[1][補足]
  {\begin{notebox}{\underline{#1}}}
  {\end{notebox}}

% メモ（デフォルト見出し：メモ）
\newenvironment{mynote}[1][メモ]
  {\begin{notebox}{\underline{#1}}}
  {\end{notebox}}

% 注意（デフォルト見出し：注意）
\newenvironment{myattention}[1][注意]
  {\begin{notebox}{\underline{#1}}}
  {\end{notebox}}



%---------- package ----------% 
%--- physics2 ---%
\usepackage{physics2}
\usephysicsmodule{ab}   % 括弧サイズの自動調整
\usephysicsmodule{ab.braket}    % braket記法
\usephysicsmodule{diagmat}  % 対角行列
\usephysicsmodule[showleft=2,showtop=2]{xmat}   % n×m行列
%--- physics2 END ---%



%---------- 追加: 確率統計で便利なマクロ ----------%
\newcommand{\PP}{\mathbb{P}}
\newcommand{\EE}{\mathbb{E}}
\newcommand{\Var}{\operatorname{Var}}
\newcommand{\Cov}{\operatorname{Cov}}
\newcommand{\ind}{\mathds{1}} % 指示関数
\newcommand{\R}{\mathbb{R}}
\newcommand{\indep}{\perp\!\!\!\perp}

% ベクトル
\usepackage{bm}
%--- 微分演算子 ---%
\usepackage{fixdif}
\usepackage{derivative}
%--- 微分演算子 END ---%
% 花文字
\usepackage{mathrsfs}
%---------- package END ----------% 

%---------- newcommand ----------% 
% 積分値の評価
\newcommand{\eval}[1]{\left.#1\right|}
% Landau記号
\newcommand{\order}[1]{\mathcal{O}\ab(#1)}
%--- physicsパッケージの\vbコマンドを再現 --%
\makeatletter
\newcommand\vb{\@ifstar\boldsymbol\mathbf}
\newcommand\va[1]{\@ifstar{\vec{#1}}{\vec{\mathrm{#1}}}}
\newcommand\vu[1]{%
\@ifstar{\hat{\boldsymbol{#1}}}{\hat{\mathbf{#1}}}}
\makeatother
%--- physicsパッケージの\vbコマンドを再現 END --%
%--- 勾配・発散・回転 ---%
% 勾配
\DeclareMathOperator{\grad}{\nabla}
% 発散
% \divが「÷」と競合するため再定義
\DeclareMathOperator{\divergence}{\nabla\cdot}
\let\divisionsymbol\div 
\renewcommand{\div}{\divergence} 
% 回転
\DeclareMathOperator{\rot}{\nabla\times}
%--- 勾配・発散・回転 END ---%
% 実部
\renewcommand{\Re}{\operatorname{Re}}
% 虚部
\renewcommand{\Im}{\operatorname{Im}}
% トレース
\newcommand{\Tr}{\operatorname{Tr}}
% rank
\newcommand{\rank}{\operatorname{rank}}
% \mqtyコマンド
\newcommand{\mqty}[1]{\begin{matrix}#1\end{matrix}}
%---------- newcommand END ----------% 

\usepackage[top=30truemm,bottom=30truemm,left=25truemm,right=25truemm]{geometry}
\usepackage{footmisc}
\renewcommand{\thefootnote}{[\arabic{footnote}] }
\renewcommand{\footnoterule}{
    \vspace{2mm}%
    \noindent\rule{\textwidth}{0.4pt}%
    \vspace{2mm}%
}
\setlength{\skip\footins}{12pt} 
\setlength{\footnotesep}{15pt}  
\interfootnotelinepenalty=10000

\usepackage{subfiles}